% homework/y5m05/y5m05.tex
% Year 5 Maths (Mr. David) - Unit 5 worksheet: Below Zero in Words.
% Goes with the deck content/y5-math/U05_L3 (Unit 5, lesson 3).
%
% The first worksheet built for a PRA teacher other than Mr Bowen. Same house
% style as the Year 7 packets (homework/hw-style.tex), with three differences
% for a Year 5 class:
%   * shorter: four pages for the student, one key page for the teacher;
%   * the language of the question is still the spine (Word help box on the
%     page where the words are used), but the numbers are Stage 5: one of the
%     two numbers is negative (framework 5Ni.01), never both;
%   * Mr. David in the word problems instead of Mr Bowen. No student names.
%
% Every item is new. The question TYPES follow the Unit 5 end-of-unit test and
% the Stage 5 scheme of work (Unit 5.4 Topic 1); none of the test's items or
% numbers is used. The next Unit 5 topics (shapes for unknown numbers,
% decimals) are left to the problem bank, not previewed here.
%
% Pure ASCII source, like every packet.
% Build (from this folder, twice each):
%   pdflatex y5m05.tex
%   pdflatex -jobname=y5m05-teacher "\def\TEACHER{}% homework/y5m05/y5m05.tex
% Year 5 Maths (Mr. David) - Unit 5 worksheet: Below Zero in Words.
% Goes with the deck content/y5-math/U05_L3 (Unit 5, lesson 3).
%
% The first worksheet built for a PRA teacher other than Mr Bowen. Same house
% style as the Year 7 packets (homework/hw-style.tex), with three differences
% for a Year 5 class:
%   * shorter: four pages for the student, one key page for the teacher;
%   * the language of the question is still the spine (Word help box on the
%     page where the words are used), but the numbers are Stage 5: one of the
%     two numbers is negative (framework 5Ni.01), never both;
%   * Mr. David in the word problems instead of Mr Bowen. No student names.
%
% Every item is new. The question TYPES follow the Unit 5 end-of-unit test and
% the Stage 5 scheme of work (Unit 5.4 Topic 1); none of the test's items or
% numbers is used. The next Unit 5 topics (shapes for unknown numbers,
% decimals) are left to the problem bank, not previewed here.
%
% Pure ASCII source, like every packet.
% Build (from this folder, twice each):
%   pdflatex y5m05.tex
%   pdflatex -jobname=y5m05-teacher "\def\TEACHER{}% homework/y5m05/y5m05.tex
% Year 5 Maths (Mr. David) - Unit 5 worksheet: Below Zero in Words.
% Goes with the deck content/y5-math/U05_L3 (Unit 5, lesson 3).
%
% The first worksheet built for a PRA teacher other than Mr Bowen. Same house
% style as the Year 7 packets (homework/hw-style.tex), with three differences
% for a Year 5 class:
%   * shorter: four pages for the student, one key page for the teacher;
%   * the language of the question is still the spine (Word help box on the
%     page where the words are used), but the numbers are Stage 5: one of the
%     two numbers is negative (framework 5Ni.01), never both;
%   * Mr. David in the word problems instead of Mr Bowen. No student names.
%
% Every item is new. The question TYPES follow the Unit 5 end-of-unit test and
% the Stage 5 scheme of work (Unit 5.4 Topic 1); none of the test's items or
% numbers is used. The next Unit 5 topics (shapes for unknown numbers,
% decimals) are left to the problem bank, not previewed here.
%
% Pure ASCII source, like every packet.
% Build (from this folder, twice each):
%   pdflatex y5m05.tex
%   pdflatex -jobname=y5m05-teacher "\def\TEACHER{}% homework/y5m05/y5m05.tex
% Year 5 Maths (Mr. David) - Unit 5 worksheet: Below Zero in Words.
% Goes with the deck content/y5-math/U05_L3 (Unit 5, lesson 3).
%
% The first worksheet built for a PRA teacher other than Mr Bowen. Same house
% style as the Year 7 packets (homework/hw-style.tex), with three differences
% for a Year 5 class:
%   * shorter: four pages for the student, one key page for the teacher;
%   * the language of the question is still the spine (Word help box on the
%     page where the words are used), but the numbers are Stage 5: one of the
%     two numbers is negative (framework 5Ni.01), never both;
%   * Mr. David in the word problems instead of Mr Bowen. No student names.
%
% Every item is new. The question TYPES follow the Unit 5 end-of-unit test and
% the Stage 5 scheme of work (Unit 5.4 Topic 1); none of the test's items or
% numbers is used. The next Unit 5 topics (shapes for unknown numbers,
% decimals) are left to the problem bank, not previewed here.
%
% Pure ASCII source, like every packet.
% Build (from this folder, twice each):
%   pdflatex y5m05.tex
%   pdflatex -jobname=y5m05-teacher "\def\TEACHER{}\input{y5m05.tex}"
% or: pwsh .claude/skills/new-homework/scripts/build-hw.ps1 y5m05

\documentclass[12pt,a4paper]{article}

\newif\ifkey
\ifdefined\TEACHER\keytrue\else
  \keyfalse
\fi

\input{../hw-style.tex}

\fancyhead[L]{\small\color{black!55}Year 5 Maths\sep Unit 5\sep Below Zero in Words}

% A number line from #1 to #2 with every whole number labelled, zero bold.
\newcommand{\numline}[2]{%
  \begin{tikzpicture}[x=0.92cm,y=1cm]
    \draw[line width=1.1pt,cbteal,{Latex[length=3mm]}-{Latex[length=3mm]}] (#1-0.7,0) -- (#2+0.7,0);
    \foreach \n in {#1,...,#2}{
      \draw[line width=0.9pt,cbteal] (\n,-0.18) -- (\n,0.18);
      \node[below=4pt,font=\small] at (\n,-0.18) {$\n$};
    }
    \draw[line width=2pt,cborange] (0,-0.28) -- (0,0.28);
  \end{tikzpicture}}

\begin{document}
\thispagestyle{fancy}

\begin{tcolorbox}[enhanced, colback=cbpurple!7, colframe=cbpurple!7, arc=7pt,
                  borderline west={6pt}{0pt}{cbpurple},
                  left=18pt, right=14pt, top=14pt, bottom=14pt]
  {\color{cbpurple}\bfseries YEAR 5 MATHS\sep UNIT 5\sep WORKSHEET}\\[6pt]
  {\color{cbpurple}\Huge\bfseries Below Zero in Words}\\[10pt]
  {\color{black!70}Adding and subtracting across zero: temperatures, floors and money}\\
  {\color{black!70}Mr. David}
\end{tcolorbox}

\vspace{8pt}
\noindent
\begin{tabularx}{\textwidth}{@{}X X@{}}
  \textbf{Name:} \blank[5.6cm] & \textbf{Class:} \blank[5.6cm] \\
  \textbf{Date given:} \blank[4.6cm] & \textbf{Date due:} \blank[5.0cm] \\
\end{tabularx}

\vspace{6pt}
\begin{activity}[How to do this worksheet]
  \begin{itemize}
    \item Write your answers \textbf{on this sheet}.
    \item In a word problem, \key{write the number sentence first}, then the answer.
    \item Stuck? Draw a number line with \key{0 in the middle}.
    \item A new word? Look for the blue \eng{Word help} box.
  \end{itemize}
\end{activity}

\begin{remember}
  \key{Zero is a bridge.} Jump to 0 first. Then jump the rest.\\[3pt]
  $3-8$: \quad 3 to 0 is 3 jumps, then 5 more jumps. \quad $3-8=-5$
\end{remember}

% ===========================================================================
\hwsection{Section A\sep Number lines}

\qhead{A1}{Count on and count back}
\noindent Write the missing numbers. The steps are always the same size.

\vspace{4pt}
\begin{enumerate}[label=(\alph*), itemsep=8pt]
  \item $9,\ 6,\ 3,\ \blank[1.4cm],\ \blank[1.4cm],\ \blank[1.4cm]$
  \item $-11,\ -7,\ -3,\ \blank[1.4cm],\ \blank[1.4cm],\ \blank[1.4cm]$
  \item $14,\ 9,\ 4,\ \blank[1.4cm],\ \blank[1.4cm],\ \blank[1.4cm]$
  \item $\blank[1.4cm],\ -15,\ -5,\ 5,\ \blank[1.4cm],\ \blank[1.4cm]$
\end{enumerate}

\qhead{A2}{Show the jumps}
\noindent Draw the jumps on the number line. Then write the answer.

\vspace{6pt}
\noindent\textbf{(a)} $3-7=\blank[2cm]$

\vspace{4pt}
\noindent\resizebox{0.92\textwidth}{!}{\numline{-6}{6}}

\vspace{14pt}
\noindent\textbf{(b)} $-5+9=\blank[2cm]$

\vspace{4pt}
\noindent\resizebox{0.92\textwidth}{!}{\numline{-6}{6}}

\qhead{A3}{Work them out}
\noindent The numbers get bigger as you go. Use zero as a bridge.

\vspace{8pt}
\begin{multicols}{3}
\begin{enumerate}[label=(\alph*), itemsep=14pt]
  \item $2-6 = \blank[1.5cm]$
  \item $1-9 = \blank[1.5cm]$
  \item $-4+7 = \blank[1.5cm]$
  \item $-8+3 = \blank[1.5cm]$
  \item $5-12 = \blank[1.5cm]$
  \item $-10+16 = \blank[1.5cm]$
  \item $20-35 = \blank[1.5cm]$
  \item $-25+40 = \blank[1.5cm]$
  \item $7-100 = \blank[1.5cm]$
  \item $-30+120 = \blank[1.5cm]$
  \item $15-204 = \blank[1.5cm]$
  \item $-60+1005 = \blank[1.2cm]$
\end{enumerate}
\end{multicols}

% ===========================================================================
\hwsection{Section B\sep The English of the question}

\begin{wordhelp}[Word help --- these words tell you to add or subtract]
  \begin{tabularx}{\linewidth}{@{}lX@{}}
    \eng{rises}, \eng{goes up}, \eng{gets warmer} & add $(+)$ \\
    \eng{falls}, \eng{drops}, \eng{goes down}, \eng{colder} & subtract $(-)$ \\
    \eng{deeper} & further down: subtract $(-)$ \\
    \eng{ground floor} (G), \eng{B1}, \eng{B2} & floor 0, then $-1$, $-2$ under the ground \\
    \eng{owe} & you must pay money back. Money you owe is \textbf{negative}. \\
    \eng{subtract 5 from 2} & start at 2: $2-5$. The number after \eng{from} comes first. \\
    \eng{the difference between} & the gap between two numbers. It is never negative. \\
  \end{tabularx}
\end{wordhelp}

\qhead{B2}{Number sentence first}
\noindent Write the number sentence. Then write the answer.

\vspace{4pt}
\begin{enumerate}[label=\arabic*., itemsep=9pt]
  \item The temperature is $2\dC$. It falls by $9\dC$. \\[4pt]
        Number sentence: \blank[5.6cm] \quad Answer: \blank[2.6cm]
  \item The lift is on floor 3. It goes down 7 floors. \\[4pt]
        Number sentence: \blank[5.6cm] \quad Floor: \blank[2.6cm]
  \item The lift is on floor B5. It goes up 8 floors. \\[4pt]
        Number sentence: \blank[5.6cm] \quad Floor: \blank[2.6cm]
  \item Mr. David has 15\,000 dong. He buys a book for 40\,000 dong. He borrows the
        rest. How much does he owe? \\[4pt]
        Number sentence: \blank[5.6cm] \quad Answer: \blank[2.6cm]
  \item Subtract 12 from 5. \\[4pt]
        Number sentence: \blank[5.6cm] \quad Answer: \blank[2.6cm]
\end{enumerate}

\qhead{B3}{The difference between}
\noindent Count the gap. A number line helps.

\vspace{4pt}
\begin{enumerate}[label=(\alph*), itemsep=10pt]
  \item The difference between $3$ and $10$. \hfill \blank[2.6cm]
  \item The difference between $-2$ and $6$. \hfill \blank[2.6cm]
  \item The difference between $-9$ and $4$. \hfill \blank[2.6cm]
  \item How much warmer is $8\dC$ than $-5\dC$? \hfill \blank[2.6cm]
  \item How many floors is it from floor B3 up to floor 6? \hfill \blank[2.6cm]
\end{enumerate}

% ===========================================================================
\hwsection{Section C\sep Read the table}

\qhead{C1}{One morning in January}
\noindent The table shows the temperature at 7 a.m. in five places on the same day.
{\small\color{black!60}(The numbers are made up for this worksheet.)}

\vspace{8pt}
\begin{center}
\begin{tabular}{|p{5.2cm}|p{3.4cm}|}
  \hline
  \rule{0pt}{1.5em}\textbf{Place} & \textbf{Temperature} \\ \hline
  \rule{0pt}{1.5em}M\d{\^o}c Ch\^au, Vietnam & $4\dC$ \\ \hline
  \rule{0pt}{1.5em}H\d{a} Long, Vietnam & $13\dC$ \\ \hline
  \rule{0pt}{1.5em}Beijing, China & $-6\dC$ \\ \hline
  \rule{0pt}{1.5em}Ulaanbaatar, Mongolia & $-24\dC$ \\ \hline
  \rule{0pt}{1.5em}Sydney, Australia & $26\dC$ \\ \hline
\end{tabular}
\end{center}

\begin{enumerate}[label=(\alph*), itemsep=11pt]
  \item Which place is the coldest? \hfill \blank[4.2cm]
  \item Beijing gets $9\dC$ warmer. What is the new temperature? \hfill \blank[2.4cm]
  \item In M\d{\^o}c Ch\^au it falls by $7\dC$. What is the new temperature? \hfill \blank[2.4cm]
  \item The difference between Sydney and Beijing? \hfill \blank[2.4cm]
  \item Write the five temperatures in order, \textbf{coldest first}. \\[6pt]
        \blank[2.2cm],\quad \blank[2.2cm],\quad \blank[2.2cm],\quad \blank[2.2cm],\quad \blank[2.2cm]
\end{enumerate}

\qhead{C2}{Estimate first}
\noindent Use easy numbers to estimate. Then work out the exact answer.

\vspace{6pt}
\noindent
\begin{tabularx}{\textwidth}{|p{3.6cm}|X|X|}
  \hline
  \rule{0pt}{1.6em}\textbf{Calculation} & \textbf{Estimate} & \textbf{Exact answer} \\ \hline
  \rule{0pt}{2.2em}$-49+151$ & & \\ \hline
  \rule{0pt}{2.2em}$302-498$ & & \\ \hline
  \rule{0pt}{2.2em}$-99+1001$ & & \\ \hline
\end{tabularx}

% ===========================================================================
\hwsection{Section D\sep Think about it}

\qhead{D1}{Mr. David's homework}
\noindent Mr. David got all three wrong. Write the correct answer. Then say what he did wrong.

\vspace{4pt}
\begin{enumerate}[label=(\alph*), itemsep=6pt]
  \item $6-10=4$ \hfill Correct answer: \blank[2.2cm]
        \wlines{1}
  \item $-3+8=-11$ \hfill Correct answer: \blank[2.2cm]
        \wlines{1}
  \item The difference between $-5$ and $3$ is $2$. \hfill Correct answer: \blank[2.2cm]
        \wlines{1}
\end{enumerate}

\qhead{D2}{Make up the story}
\noindent Write a word problem for each number sentence. Use a temperature, a lift
or money. Use a word from the Word help box. End with a question.

\vspace{6pt}
\noindent\textbf{(a)} \expr{4-9=-5}
\wlines{2}

\vspace{8pt}
\noindent\textbf{(b)} \expr{-3+10=7}
\wlines{2}

% ===========================================================================
% ANSWER KEY (teacher copy only)
% ===========================================================================
\ifkey
\newpage
\hwsection[cbred]{Answer Key --- Teacher Copy}
\linespread{1.0}\selectfont
\setlist[enumerate]{itemsep=1pt, topsep=2pt, leftmargin=*}
\small

\qhead[cbred]{A1}{Count on and count back}
\noindent (a) $0, -3, -6$ \quad (b) $1, 5, 9$ \quad (c) $-1, -6, -11$ \quad
(d) $-25$ \ldots\ $15, 25$ (steps of 10). A student who writes $-1$ after $3$ in (a)
has stopped seeing the step size at zero: count the jumps aloud with them.

\qhead[cbred]{A2}{Show the jumps}
\noindent (a) $-4$: 3 jumps left to 0, then 4 more. \quad (b) $4$: 5 jumps right to 0,
then 4 more. Check the arrow starts at the first number, not at 0.

\qhead[cbred]{A3}{Work them out}
\noindent (a) $-4$ (b) $-8$ (c) $3$ (d) $-5$ (e) $-7$ (f) $6$ (g) $-15$ (h) $15$
(i) $-93$ (j) $90$ (k) $-189$ (l) $945$. \\
The common wrong answers are the gap with the wrong sign ($7-100=93$) and the
``keep the minus'' sum ($-4+7=-11$, the scheme of work's misconception). Both mean
the student is not starting at the first number.

\qhead[cbred]{B2}{Number sentence first}
\noindent 1. $2-9=-7\dC$ \quad 2. $3-7=-4$, floor B4 \quad
3. $-5+8=3$, floor 3 \quad 4. $15\,000-40\,000=-25\,000$: he owes 25\,000 dong \quad
5. $5-12=-7$. \\
In 5, $12-5=7$ is the English mistake (the number after \textit{from} comes first),
not an arithmetic one: ask the student to read the sentence aloud.

\qhead[cbred]{B3}{The difference between}
\noindent (a) 7 \quad (b) 8 \quad (c) 13 \quad (d) $13\dC$ \quad (e) 9 floors (3 up to
G, then 6). Answers of 4, 5 or 3 in (b)--(e) come from subtracting the digits and
ignoring the minus sign. A difference is never negative: $-13$ is wrong.

\qhead[cbred]{C1}{One morning in January}
\noindent (a) Ulaanbaatar \quad (b) $-6+9=3\dC$ \quad (c) $4-7=-3\dC$ \quad
(d) $26+6=32\dC$ \quad (e) $-24, -6, 4, 13, 26$ ($\dC$). In (e), students who put $-6$
before $-24$ are ordering by the digits; ask which is colder.

\qhead[cbred]{C2}{Estimate first}
\noindent $-50+150=100$, exact $102$. \quad $300-500=-200$, exact $-196$. \quad
$-100+1000=900$, exact $902$. Any sensible rounded estimate is fine; the point is that
the exact answer is close to it.

\qhead[cbred]{D1}{Mr. David's homework}
\noindent (a) $-4$: he found the gap between 6 and 10 instead of starting at 6 and going
down 10. (b) $5$: he added 3 and 8 and kept the minus sign. (c) $8$: he subtracted 3
from 5; the gap crosses zero, $5+3$.

\qhead[cbred]{D2}{Make up the story}
\noindent Check in this order: (1) it \textbf{starts} at the first number (4, or $-3$);
(2) the word matches the sign (falls / goes down / spends for (a), rises / goes up /
earns for (b)); (3) full sentences ending in a question; (4) the answer is $-5$ or 7.
Example (a): ``It is $4\dC$. The temperature falls by $9\dC$. What is the temperature
now?'' Example (b): ``The lift is on B3. It goes up 10 floors. Which floor is it on now?''

\fi
\end{document}
"
% or: pwsh .claude/skills/new-homework/scripts/build-hw.ps1 y5m05

\documentclass[12pt,a4paper]{article}

\newif\ifkey
\ifdefined\TEACHER\keytrue\else
  \keyfalse
\fi

% homework/hw-style.tex
% Shared house style for every Year 7 homework packet. \input this from the
% packet file, right after \documentclass and the answer-key switch.
%
% Do not fork this file per packet. If a packet needs a different look, the
% question is whether every packet should have it.
%
% The choices here are deliberate and were arrived at by printing and looking:
%   * No solid dark banners. Every header is a light tint with a coloured spine
%     on the left, so colour reads clearly without flooding a school printer.
%   * No marks and no mark totals. These are practice, not exam papers.
%   * Lato at 12pt with open leading, plus mathastext so inline maths is Lato
%     too. Without mathastext every "-6 + 4" is Computer Modern serif and the
%     sheet looks half-typeset.
%   * Writing lines are spaced for Year 7 handwriting, not adult margin notes.
%   * Pages fill. \needspace is redefined below so a heading no longer drags
%     the page break up to itself, and every heading is glued to the box under
%     it. Packets built before 2026-09-30 were laid out under the old rules and
%     will reflow if rebuilt -- re-read every page if you rebuild one.
%
% Provides: \wlines \blank \tickbox \sep \qmk-free marks (none) \hwsection
% \qhead \expr, the workspace box, and the remember / wordhelp / activity /
% yourturn coloured boxes. Palette matches the lesson decks exactly.

% -- Packages ----------------------------------------------------------------
\usepackage[T1]{fontenc}
\usepackage[utf8]{inputenc}
\usepackage[default]{lato}
\usepackage[a4paper,top=1.7cm,bottom=1.7cm,left=2.0cm,right=2.0cm,headheight=15pt]{geometry}
\usepackage{amsmath,amssymb}
\usepackage{xcolor}
\usepackage[most]{tcolorbox}
\usepackage{enumitem}
\usepackage{tikz}
\usetikzlibrary{arrows.meta}
\usepackage{array}
\usepackage{tabularx}
\usepackage{fancyhdr}
\usepackage{multicol}
\usepackage{needspace}
\usepackage{graphicx}

% Make maths use Lato too. Without this, every "-6 + 4" on the page is set in
% Computer Modern serif and the sheet looks half-typeset. Must come last.
\usepackage[italic]{mathastext}

% -- \needspace, without the hole it used to leave ----------------------------
% needspace.sty writes "\vskip 0pt plus X \penalty-100" in front of every
% heading. A page that ends anywhere else is underfull with no stretch, so TeX
% scores it b=10000; a page that ends AT a heading gets that X of stretch and a
% finite badness, and wins. So any heading in roughly the bottom half of a page
% claimed the break, and the question under it moved over whole, leaving a hole
% the size of the question. \tracingpages showed it exactly (b=4531 at the
% heading, b=10000 at every later break) while building homework/q1rev, and it
% is where the half-empty pages in HW 1-8 came from.
% This keeps only the \penalty9999, which still refuses to put a heading in the
% last X of a page. Keeping a heading WITH the box under it is \nobreak's job
% now -- see \hwsection and \qhead below.
\makeatletter
\renewcommand{\needspace}[1]{\begingroup\setlength{\dimen@}{#1}%
  \vskip\dimen@\penalty9999\vskip-\dimen@\endgroup}
\makeatother

\linespread{1.06}\selectfont

% -- The Learner's Book palette (same hex values as the lesson decks) ---------
\definecolor{cbteal}{HTML}{0087A8}
\definecolor{cbpurple}{HTML}{5C2483}
\definecolor{cborange}{HTML}{C25E12}
\definecolor{cbgreen}{HTML}{4A8B23}
\definecolor{cbred}{HTML}{C8102E}
\definecolor{cbblue}{HTML}{1A5FA8}
\definecolor{rulegray}{HTML}{9AA5AD}

% -- Answer lines and blanks -------------------------------------------------
% \wlines{n} draws n ruled writing lines, spaced wide enough for Year 7
% handwriting rather than for an adult with a fine pen.
\makeatletter
\newcommand{\wlines}[1]{%
  \par\vspace{0.75em}%
  \@tempcnta=0\relax
  \@whilenum\@tempcnta<#1\do{%
    \hbox to \linewidth{\color{rulegray}\leaders\hrule height 0.5pt\hfill}%
    \vspace{1.5em}%
    \advance\@tempcnta by 1\relax}%
  \par\vspace{-0.8em}%
}
\makeatother

% \blank[width] is an inline gap to write a short answer in.
\newcommand{\blank}[1][2.6cm]{\underline{\hspace{#1}}}

% A boxed space for showing working.
\newtcolorbox{workspace}[1][3.4cm]{%
  colback=white, colframe=rulegray, boxrule=0.5pt, arc=5pt,
  height=#1, left=6pt, right=6pt, top=4pt, bottom=4pt,
  before skip=9pt, after skip=11pt}

% An empty tick box.
\newcommand{\tickbox}{\fbox{\rule{0pt}{1.15em}\rule{1.15em}{0pt}}}

% A middot separator.
\newcommand{\sep}{\,\textperiodcentered\,}

% -- Section and question headings -------------------------------------------
% Light tint plus a fat coloured spine on the left. No solid dark banner: the
% colour reads clearly but barely costs the printer anything.
% needspace keeps a heading out of the last few lines of a page. The \nobreak
% after it keeps it with whatever comes next: a remember or word-help box that
% does not fit takes its heading over with it, instead of leaving the heading
% alone at the foot of the page. (The old needspace did that by accident, and
% left a hole every time; see above.)
\newcommand{\hwsection}[2][cbteal]{%
  \needspace{8\baselineskip}%
  \par\vspace{13pt}%
  \noindent
  \begin{tcolorbox}[enhanced, nobeforeafter, width=\textwidth,
      colback=#1!7, colframe=#1!7, arc=5pt,
      borderline west={5pt}{0pt}{#1},
      left=14pt, right=10pt, top=7pt, bottom=7pt]
    {\large\bfseries\color{#1}#2}
  \end{tcolorbox}%
  \par\nobreak\vspace{9pt}}

% \qhead[colour]{A2}{Moving along the line} -- a question heading.
\newcommand{\qhead}[3][cbteal]{%
  \needspace{6\baselineskip}%
  \par\vspace{11pt}%
  \noindent{\large\bfseries\color{#1}#2}\quad{\large\bfseries #3}\par\nobreak\vspace{4pt}}

% -- Coloured boxes ----------------------------------------------------------
% Shared look: rounded, light fill, coloured rule, coloured title text on a
% pale title strip rather than a dark one.
% NOT breakable. A breakable box that cannot fit its first lines forces its own
% page break, and \nobreak cannot stop a forced break -- so the heading above
% it was left alone at the foot of the page (tested: it happens with \nobreak
% in place). Unbreakable, the page builder moves heading and box together.
% These boxes are a few lines each; none needs to split.
\tcbset{hwbox/.style={
  enhanced, arc=6pt, boxrule=1pt,
  left=11pt, right=11pt, top=8pt, bottom=8pt,
  fonttitle=\bfseries, attach boxed title to top left={xshift=12pt, yshift=-8pt},
  boxed title style={arc=4pt, boxrule=1pt, left=7pt, right=7pt, top=2pt, bottom=2pt},
  top=13pt}}

% Orange: things to remember. Matches the deck's "Write This Down".
\newtcolorbox{remember}[1][Remember from the lesson]{hwbox,
  colback=cborange!6, colframe=cborange, coltitle=cborange,
  colbacktitle=white, title={#1}}

% Blue: English help. The barrier in this class is the wording, not the maths.
\newtcolorbox{wordhelp}[1][Word help]{hwbox,
  colback=cbblue!6, colframe=cbblue, coltitle=cbblue,
  colbacktitle=white, title={#1}}

% Purple: an activity or instruction.
\newtcolorbox{activity}[1][How to do this homework]{hwbox,
  colback=cbpurple!6, colframe=cbpurple, coltitle=cbpurple,
  colbacktitle=white, title={#1}}

% Green: your turn / a friendly nudge.
\newtcolorbox{yourturn}[1][Your turn]{hwbox,
  colback=cbgreen!7, colframe=cbgreen, coltitle=cbgreen!75!black,
  colbacktitle=white, title={#1}}

% -- Shorthands --------------------------------------------------------------
\newcommand{\dC}{\,^{\circ}\mathrm{C}}          % use inside maths: $-8\dC$
\newcommand{\key}[1]{{\color{cborange}\textbf{#1}}}
\newcommand{\eng}[1]{{\color{cbblue}\textbf{#1}}}

\setlist[enumerate]{itemsep=8pt, topsep=5pt, leftmargin=*}
\setlist[itemize]{itemsep=4pt, topsep=4pt, leftmargin=*}
\renewcommand{\arraystretch}{1.45}

% -- An expression the student is asked to work from -------------------------
\newcommand{\expr}[1]{%
  \tcbox[on line, colback=cbgreen!10, colframe=cbgreen!55, boxrule=1pt, arc=4pt,
         left=9pt, right=9pt, top=4pt, bottom=4pt]{\large\bfseries $#1$}}

% -- Running head and foot ---------------------------------------------------
% The packet overrides \fancyhead[L] with its own title.
\pagestyle{fancy}
\fancyhf{}
\fancyhead[L]{\small\color{black!55}Year 7\sep Homework}
\fancyhead[R]{\small\color{black!55}Name: \blank[3.4cm]}
\fancyfoot[C]{\small\color{black!55}Page \thepage}
\renewcommand{\headrulewidth}{0pt}
\renewcommand{\headrule}{\vspace{2pt}\hbox to\headwidth{\color{cbteal!45}\leaders\hrule height 1.2pt\hfill}}



\fancyhead[L]{\small\color{black!55}Year 5 Maths\sep Unit 5\sep Below Zero in Words}

% A number line from #1 to #2 with every whole number labelled, zero bold.
\newcommand{\numline}[2]{%
  \begin{tikzpicture}[x=0.92cm,y=1cm]
    \draw[line width=1.1pt,cbteal,{Latex[length=3mm]}-{Latex[length=3mm]}] (#1-0.7,0) -- (#2+0.7,0);
    \foreach \n in {#1,...,#2}{
      \draw[line width=0.9pt,cbteal] (\n,-0.18) -- (\n,0.18);
      \node[below=4pt,font=\small] at (\n,-0.18) {$\n$};
    }
    \draw[line width=2pt,cborange] (0,-0.28) -- (0,0.28);
  \end{tikzpicture}}

\begin{document}
\thispagestyle{fancy}

\begin{tcolorbox}[enhanced, colback=cbpurple!7, colframe=cbpurple!7, arc=7pt,
                  borderline west={6pt}{0pt}{cbpurple},
                  left=18pt, right=14pt, top=14pt, bottom=14pt]
  {\color{cbpurple}\bfseries YEAR 5 MATHS\sep UNIT 5\sep WORKSHEET}\\[6pt]
  {\color{cbpurple}\Huge\bfseries Below Zero in Words}\\[10pt]
  {\color{black!70}Adding and subtracting across zero: temperatures, floors and money}\\
  {\color{black!70}Mr. David}
\end{tcolorbox}

\vspace{8pt}
\noindent
\begin{tabularx}{\textwidth}{@{}X X@{}}
  \textbf{Name:} \blank[5.6cm] & \textbf{Class:} \blank[5.6cm] \\
  \textbf{Date given:} \blank[4.6cm] & \textbf{Date due:} \blank[5.0cm] \\
\end{tabularx}

\vspace{6pt}
\begin{activity}[How to do this worksheet]
  \begin{itemize}
    \item Write your answers \textbf{on this sheet}.
    \item In a word problem, \key{write the number sentence first}, then the answer.
    \item Stuck? Draw a number line with \key{0 in the middle}.
    \item A new word? Look for the blue \eng{Word help} box.
  \end{itemize}
\end{activity}

\begin{remember}
  \key{Zero is a bridge.} Jump to 0 first. Then jump the rest.\\[3pt]
  $3-8$: \quad 3 to 0 is 3 jumps, then 5 more jumps. \quad $3-8=-5$
\end{remember}

% ===========================================================================
\hwsection{Section A\sep Number lines}

\qhead{A1}{Count on and count back}
\noindent Write the missing numbers. The steps are always the same size.

\vspace{4pt}
\begin{enumerate}[label=(\alph*), itemsep=8pt]
  \item $9,\ 6,\ 3,\ \blank[1.4cm],\ \blank[1.4cm],\ \blank[1.4cm]$
  \item $-11,\ -7,\ -3,\ \blank[1.4cm],\ \blank[1.4cm],\ \blank[1.4cm]$
  \item $14,\ 9,\ 4,\ \blank[1.4cm],\ \blank[1.4cm],\ \blank[1.4cm]$
  \item $\blank[1.4cm],\ -15,\ -5,\ 5,\ \blank[1.4cm],\ \blank[1.4cm]$
\end{enumerate}

\qhead{A2}{Show the jumps}
\noindent Draw the jumps on the number line. Then write the answer.

\vspace{6pt}
\noindent\textbf{(a)} $3-7=\blank[2cm]$

\vspace{4pt}
\noindent\resizebox{0.92\textwidth}{!}{\numline{-6}{6}}

\vspace{14pt}
\noindent\textbf{(b)} $-5+9=\blank[2cm]$

\vspace{4pt}
\noindent\resizebox{0.92\textwidth}{!}{\numline{-6}{6}}

\qhead{A3}{Work them out}
\noindent The numbers get bigger as you go. Use zero as a bridge.

\vspace{8pt}
\begin{multicols}{3}
\begin{enumerate}[label=(\alph*), itemsep=14pt]
  \item $2-6 = \blank[1.5cm]$
  \item $1-9 = \blank[1.5cm]$
  \item $-4+7 = \blank[1.5cm]$
  \item $-8+3 = \blank[1.5cm]$
  \item $5-12 = \blank[1.5cm]$
  \item $-10+16 = \blank[1.5cm]$
  \item $20-35 = \blank[1.5cm]$
  \item $-25+40 = \blank[1.5cm]$
  \item $7-100 = \blank[1.5cm]$
  \item $-30+120 = \blank[1.5cm]$
  \item $15-204 = \blank[1.5cm]$
  \item $-60+1005 = \blank[1.2cm]$
\end{enumerate}
\end{multicols}

% ===========================================================================
\hwsection{Section B\sep The English of the question}

\begin{wordhelp}[Word help --- these words tell you to add or subtract]
  \begin{tabularx}{\linewidth}{@{}lX@{}}
    \eng{rises}, \eng{goes up}, \eng{gets warmer} & add $(+)$ \\
    \eng{falls}, \eng{drops}, \eng{goes down}, \eng{colder} & subtract $(-)$ \\
    \eng{deeper} & further down: subtract $(-)$ \\
    \eng{ground floor} (G), \eng{B1}, \eng{B2} & floor 0, then $-1$, $-2$ under the ground \\
    \eng{owe} & you must pay money back. Money you owe is \textbf{negative}. \\
    \eng{subtract 5 from 2} & start at 2: $2-5$. The number after \eng{from} comes first. \\
    \eng{the difference between} & the gap between two numbers. It is never negative. \\
  \end{tabularx}
\end{wordhelp}

\qhead{B2}{Number sentence first}
\noindent Write the number sentence. Then write the answer.

\vspace{4pt}
\begin{enumerate}[label=\arabic*., itemsep=9pt]
  \item The temperature is $2\dC$. It falls by $9\dC$. \\[4pt]
        Number sentence: \blank[5.6cm] \quad Answer: \blank[2.6cm]
  \item The lift is on floor 3. It goes down 7 floors. \\[4pt]
        Number sentence: \blank[5.6cm] \quad Floor: \blank[2.6cm]
  \item The lift is on floor B5. It goes up 8 floors. \\[4pt]
        Number sentence: \blank[5.6cm] \quad Floor: \blank[2.6cm]
  \item Mr. David has 15\,000 dong. He buys a book for 40\,000 dong. He borrows the
        rest. How much does he owe? \\[4pt]
        Number sentence: \blank[5.6cm] \quad Answer: \blank[2.6cm]
  \item Subtract 12 from 5. \\[4pt]
        Number sentence: \blank[5.6cm] \quad Answer: \blank[2.6cm]
\end{enumerate}

\qhead{B3}{The difference between}
\noindent Count the gap. A number line helps.

\vspace{4pt}
\begin{enumerate}[label=(\alph*), itemsep=10pt]
  \item The difference between $3$ and $10$. \hfill \blank[2.6cm]
  \item The difference between $-2$ and $6$. \hfill \blank[2.6cm]
  \item The difference between $-9$ and $4$. \hfill \blank[2.6cm]
  \item How much warmer is $8\dC$ than $-5\dC$? \hfill \blank[2.6cm]
  \item How many floors is it from floor B3 up to floor 6? \hfill \blank[2.6cm]
\end{enumerate}

% ===========================================================================
\hwsection{Section C\sep Read the table}

\qhead{C1}{One morning in January}
\noindent The table shows the temperature at 7 a.m. in five places on the same day.
{\small\color{black!60}(The numbers are made up for this worksheet.)}

\vspace{8pt}
\begin{center}
\begin{tabular}{|p{5.2cm}|p{3.4cm}|}
  \hline
  \rule{0pt}{1.5em}\textbf{Place} & \textbf{Temperature} \\ \hline
  \rule{0pt}{1.5em}M\d{\^o}c Ch\^au, Vietnam & $4\dC$ \\ \hline
  \rule{0pt}{1.5em}H\d{a} Long, Vietnam & $13\dC$ \\ \hline
  \rule{0pt}{1.5em}Beijing, China & $-6\dC$ \\ \hline
  \rule{0pt}{1.5em}Ulaanbaatar, Mongolia & $-24\dC$ \\ \hline
  \rule{0pt}{1.5em}Sydney, Australia & $26\dC$ \\ \hline
\end{tabular}
\end{center}

\begin{enumerate}[label=(\alph*), itemsep=11pt]
  \item Which place is the coldest? \hfill \blank[4.2cm]
  \item Beijing gets $9\dC$ warmer. What is the new temperature? \hfill \blank[2.4cm]
  \item In M\d{\^o}c Ch\^au it falls by $7\dC$. What is the new temperature? \hfill \blank[2.4cm]
  \item The difference between Sydney and Beijing? \hfill \blank[2.4cm]
  \item Write the five temperatures in order, \textbf{coldest first}. \\[6pt]
        \blank[2.2cm],\quad \blank[2.2cm],\quad \blank[2.2cm],\quad \blank[2.2cm],\quad \blank[2.2cm]
\end{enumerate}

\qhead{C2}{Estimate first}
\noindent Use easy numbers to estimate. Then work out the exact answer.

\vspace{6pt}
\noindent
\begin{tabularx}{\textwidth}{|p{3.6cm}|X|X|}
  \hline
  \rule{0pt}{1.6em}\textbf{Calculation} & \textbf{Estimate} & \textbf{Exact answer} \\ \hline
  \rule{0pt}{2.2em}$-49+151$ & & \\ \hline
  \rule{0pt}{2.2em}$302-498$ & & \\ \hline
  \rule{0pt}{2.2em}$-99+1001$ & & \\ \hline
\end{tabularx}

% ===========================================================================
\hwsection{Section D\sep Think about it}

\qhead{D1}{Mr. David's homework}
\noindent Mr. David got all three wrong. Write the correct answer. Then say what he did wrong.

\vspace{4pt}
\begin{enumerate}[label=(\alph*), itemsep=6pt]
  \item $6-10=4$ \hfill Correct answer: \blank[2.2cm]
        \wlines{1}
  \item $-3+8=-11$ \hfill Correct answer: \blank[2.2cm]
        \wlines{1}
  \item The difference between $-5$ and $3$ is $2$. \hfill Correct answer: \blank[2.2cm]
        \wlines{1}
\end{enumerate}

\qhead{D2}{Make up the story}
\noindent Write a word problem for each number sentence. Use a temperature, a lift
or money. Use a word from the Word help box. End with a question.

\vspace{6pt}
\noindent\textbf{(a)} \expr{4-9=-5}
\wlines{2}

\vspace{8pt}
\noindent\textbf{(b)} \expr{-3+10=7}
\wlines{2}

% ===========================================================================
% ANSWER KEY (teacher copy only)
% ===========================================================================
\ifkey
\newpage
\hwsection[cbred]{Answer Key --- Teacher Copy}
\linespread{1.0}\selectfont
\setlist[enumerate]{itemsep=1pt, topsep=2pt, leftmargin=*}
\small

\qhead[cbred]{A1}{Count on and count back}
\noindent (a) $0, -3, -6$ \quad (b) $1, 5, 9$ \quad (c) $-1, -6, -11$ \quad
(d) $-25$ \ldots\ $15, 25$ (steps of 10). A student who writes $-1$ after $3$ in (a)
has stopped seeing the step size at zero: count the jumps aloud with them.

\qhead[cbred]{A2}{Show the jumps}
\noindent (a) $-4$: 3 jumps left to 0, then 4 more. \quad (b) $4$: 5 jumps right to 0,
then 4 more. Check the arrow starts at the first number, not at 0.

\qhead[cbred]{A3}{Work them out}
\noindent (a) $-4$ (b) $-8$ (c) $3$ (d) $-5$ (e) $-7$ (f) $6$ (g) $-15$ (h) $15$
(i) $-93$ (j) $90$ (k) $-189$ (l) $945$. \\
The common wrong answers are the gap with the wrong sign ($7-100=93$) and the
``keep the minus'' sum ($-4+7=-11$, the scheme of work's misconception). Both mean
the student is not starting at the first number.

\qhead[cbred]{B2}{Number sentence first}
\noindent 1. $2-9=-7\dC$ \quad 2. $3-7=-4$, floor B4 \quad
3. $-5+8=3$, floor 3 \quad 4. $15\,000-40\,000=-25\,000$: he owes 25\,000 dong \quad
5. $5-12=-7$. \\
In 5, $12-5=7$ is the English mistake (the number after \textit{from} comes first),
not an arithmetic one: ask the student to read the sentence aloud.

\qhead[cbred]{B3}{The difference between}
\noindent (a) 7 \quad (b) 8 \quad (c) 13 \quad (d) $13\dC$ \quad (e) 9 floors (3 up to
G, then 6). Answers of 4, 5 or 3 in (b)--(e) come from subtracting the digits and
ignoring the minus sign. A difference is never negative: $-13$ is wrong.

\qhead[cbred]{C1}{One morning in January}
\noindent (a) Ulaanbaatar \quad (b) $-6+9=3\dC$ \quad (c) $4-7=-3\dC$ \quad
(d) $26+6=32\dC$ \quad (e) $-24, -6, 4, 13, 26$ ($\dC$). In (e), students who put $-6$
before $-24$ are ordering by the digits; ask which is colder.

\qhead[cbred]{C2}{Estimate first}
\noindent $-50+150=100$, exact $102$. \quad $300-500=-200$, exact $-196$. \quad
$-100+1000=900$, exact $902$. Any sensible rounded estimate is fine; the point is that
the exact answer is close to it.

\qhead[cbred]{D1}{Mr. David's homework}
\noindent (a) $-4$: he found the gap between 6 and 10 instead of starting at 6 and going
down 10. (b) $5$: he added 3 and 8 and kept the minus sign. (c) $8$: he subtracted 3
from 5; the gap crosses zero, $5+3$.

\qhead[cbred]{D2}{Make up the story}
\noindent Check in this order: (1) it \textbf{starts} at the first number (4, or $-3$);
(2) the word matches the sign (falls / goes down / spends for (a), rises / goes up /
earns for (b)); (3) full sentences ending in a question; (4) the answer is $-5$ or 7.
Example (a): ``It is $4\dC$. The temperature falls by $9\dC$. What is the temperature
now?'' Example (b): ``The lift is on B3. It goes up 10 floors. Which floor is it on now?''

\fi
\end{document}
"
% or: pwsh .claude/skills/new-homework/scripts/build-hw.ps1 y5m05

\documentclass[12pt,a4paper]{article}

\newif\ifkey
\ifdefined\TEACHER\keytrue\else
  \keyfalse
\fi

% homework/hw-style.tex
% Shared house style for every Year 7 homework packet. \input this from the
% packet file, right after \documentclass and the answer-key switch.
%
% Do not fork this file per packet. If a packet needs a different look, the
% question is whether every packet should have it.
%
% The choices here are deliberate and were arrived at by printing and looking:
%   * No solid dark banners. Every header is a light tint with a coloured spine
%     on the left, so colour reads clearly without flooding a school printer.
%   * No marks and no mark totals. These are practice, not exam papers.
%   * Lato at 12pt with open leading, plus mathastext so inline maths is Lato
%     too. Without mathastext every "-6 + 4" is Computer Modern serif and the
%     sheet looks half-typeset.
%   * Writing lines are spaced for Year 7 handwriting, not adult margin notes.
%   * Pages fill. \needspace is redefined below so a heading no longer drags
%     the page break up to itself, and every heading is glued to the box under
%     it. Packets built before 2026-09-30 were laid out under the old rules and
%     will reflow if rebuilt -- re-read every page if you rebuild one.
%
% Provides: \wlines \blank \tickbox \sep \qmk-free marks (none) \hwsection
% \qhead \expr, the workspace box, and the remember / wordhelp / activity /
% yourturn coloured boxes. Palette matches the lesson decks exactly.

% -- Packages ----------------------------------------------------------------
\usepackage[T1]{fontenc}
\usepackage[utf8]{inputenc}
\usepackage[default]{lato}
\usepackage[a4paper,top=1.7cm,bottom=1.7cm,left=2.0cm,right=2.0cm,headheight=15pt]{geometry}
\usepackage{amsmath,amssymb}
\usepackage{xcolor}
\usepackage[most]{tcolorbox}
\usepackage{enumitem}
\usepackage{tikz}
\usetikzlibrary{arrows.meta}
\usepackage{array}
\usepackage{tabularx}
\usepackage{fancyhdr}
\usepackage{multicol}
\usepackage{needspace}
\usepackage{graphicx}

% Make maths use Lato too. Without this, every "-6 + 4" on the page is set in
% Computer Modern serif and the sheet looks half-typeset. Must come last.
\usepackage[italic]{mathastext}

% -- \needspace, without the hole it used to leave ----------------------------
% needspace.sty writes "\vskip 0pt plus X \penalty-100" in front of every
% heading. A page that ends anywhere else is underfull with no stretch, so TeX
% scores it b=10000; a page that ends AT a heading gets that X of stretch and a
% finite badness, and wins. So any heading in roughly the bottom half of a page
% claimed the break, and the question under it moved over whole, leaving a hole
% the size of the question. \tracingpages showed it exactly (b=4531 at the
% heading, b=10000 at every later break) while building homework/q1rev, and it
% is where the half-empty pages in HW 1-8 came from.
% This keeps only the \penalty9999, which still refuses to put a heading in the
% last X of a page. Keeping a heading WITH the box under it is \nobreak's job
% now -- see \hwsection and \qhead below.
\makeatletter
\renewcommand{\needspace}[1]{\begingroup\setlength{\dimen@}{#1}%
  \vskip\dimen@\penalty9999\vskip-\dimen@\endgroup}
\makeatother

\linespread{1.06}\selectfont

% -- The Learner's Book palette (same hex values as the lesson decks) ---------
\definecolor{cbteal}{HTML}{0087A8}
\definecolor{cbpurple}{HTML}{5C2483}
\definecolor{cborange}{HTML}{C25E12}
\definecolor{cbgreen}{HTML}{4A8B23}
\definecolor{cbred}{HTML}{C8102E}
\definecolor{cbblue}{HTML}{1A5FA8}
\definecolor{rulegray}{HTML}{9AA5AD}

% -- Answer lines and blanks -------------------------------------------------
% \wlines{n} draws n ruled writing lines, spaced wide enough for Year 7
% handwriting rather than for an adult with a fine pen.
\makeatletter
\newcommand{\wlines}[1]{%
  \par\vspace{0.75em}%
  \@tempcnta=0\relax
  \@whilenum\@tempcnta<#1\do{%
    \hbox to \linewidth{\color{rulegray}\leaders\hrule height 0.5pt\hfill}%
    \vspace{1.5em}%
    \advance\@tempcnta by 1\relax}%
  \par\vspace{-0.8em}%
}
\makeatother

% \blank[width] is an inline gap to write a short answer in.
\newcommand{\blank}[1][2.6cm]{\underline{\hspace{#1}}}

% A boxed space for showing working.
\newtcolorbox{workspace}[1][3.4cm]{%
  colback=white, colframe=rulegray, boxrule=0.5pt, arc=5pt,
  height=#1, left=6pt, right=6pt, top=4pt, bottom=4pt,
  before skip=9pt, after skip=11pt}

% An empty tick box.
\newcommand{\tickbox}{\fbox{\rule{0pt}{1.15em}\rule{1.15em}{0pt}}}

% A middot separator.
\newcommand{\sep}{\,\textperiodcentered\,}

% -- Section and question headings -------------------------------------------
% Light tint plus a fat coloured spine on the left. No solid dark banner: the
% colour reads clearly but barely costs the printer anything.
% needspace keeps a heading out of the last few lines of a page. The \nobreak
% after it keeps it with whatever comes next: a remember or word-help box that
% does not fit takes its heading over with it, instead of leaving the heading
% alone at the foot of the page. (The old needspace did that by accident, and
% left a hole every time; see above.)
\newcommand{\hwsection}[2][cbteal]{%
  \needspace{8\baselineskip}%
  \par\vspace{13pt}%
  \noindent
  \begin{tcolorbox}[enhanced, nobeforeafter, width=\textwidth,
      colback=#1!7, colframe=#1!7, arc=5pt,
      borderline west={5pt}{0pt}{#1},
      left=14pt, right=10pt, top=7pt, bottom=7pt]
    {\large\bfseries\color{#1}#2}
  \end{tcolorbox}%
  \par\nobreak\vspace{9pt}}

% \qhead[colour]{A2}{Moving along the line} -- a question heading.
\newcommand{\qhead}[3][cbteal]{%
  \needspace{6\baselineskip}%
  \par\vspace{11pt}%
  \noindent{\large\bfseries\color{#1}#2}\quad{\large\bfseries #3}\par\nobreak\vspace{4pt}}

% -- Coloured boxes ----------------------------------------------------------
% Shared look: rounded, light fill, coloured rule, coloured title text on a
% pale title strip rather than a dark one.
% NOT breakable. A breakable box that cannot fit its first lines forces its own
% page break, and \nobreak cannot stop a forced break -- so the heading above
% it was left alone at the foot of the page (tested: it happens with \nobreak
% in place). Unbreakable, the page builder moves heading and box together.
% These boxes are a few lines each; none needs to split.
\tcbset{hwbox/.style={
  enhanced, arc=6pt, boxrule=1pt,
  left=11pt, right=11pt, top=8pt, bottom=8pt,
  fonttitle=\bfseries, attach boxed title to top left={xshift=12pt, yshift=-8pt},
  boxed title style={arc=4pt, boxrule=1pt, left=7pt, right=7pt, top=2pt, bottom=2pt},
  top=13pt}}

% Orange: things to remember. Matches the deck's "Write This Down".
\newtcolorbox{remember}[1][Remember from the lesson]{hwbox,
  colback=cborange!6, colframe=cborange, coltitle=cborange,
  colbacktitle=white, title={#1}}

% Blue: English help. The barrier in this class is the wording, not the maths.
\newtcolorbox{wordhelp}[1][Word help]{hwbox,
  colback=cbblue!6, colframe=cbblue, coltitle=cbblue,
  colbacktitle=white, title={#1}}

% Purple: an activity or instruction.
\newtcolorbox{activity}[1][How to do this homework]{hwbox,
  colback=cbpurple!6, colframe=cbpurple, coltitle=cbpurple,
  colbacktitle=white, title={#1}}

% Green: your turn / a friendly nudge.
\newtcolorbox{yourturn}[1][Your turn]{hwbox,
  colback=cbgreen!7, colframe=cbgreen, coltitle=cbgreen!75!black,
  colbacktitle=white, title={#1}}

% -- Shorthands --------------------------------------------------------------
\newcommand{\dC}{\,^{\circ}\mathrm{C}}          % use inside maths: $-8\dC$
\newcommand{\key}[1]{{\color{cborange}\textbf{#1}}}
\newcommand{\eng}[1]{{\color{cbblue}\textbf{#1}}}

\setlist[enumerate]{itemsep=8pt, topsep=5pt, leftmargin=*}
\setlist[itemize]{itemsep=4pt, topsep=4pt, leftmargin=*}
\renewcommand{\arraystretch}{1.45}

% -- An expression the student is asked to work from -------------------------
\newcommand{\expr}[1]{%
  \tcbox[on line, colback=cbgreen!10, colframe=cbgreen!55, boxrule=1pt, arc=4pt,
         left=9pt, right=9pt, top=4pt, bottom=4pt]{\large\bfseries $#1$}}

% -- Running head and foot ---------------------------------------------------
% The packet overrides \fancyhead[L] with its own title.
\pagestyle{fancy}
\fancyhf{}
\fancyhead[L]{\small\color{black!55}Year 7\sep Homework}
\fancyhead[R]{\small\color{black!55}Name: \blank[3.4cm]}
\fancyfoot[C]{\small\color{black!55}Page \thepage}
\renewcommand{\headrulewidth}{0pt}
\renewcommand{\headrule}{\vspace{2pt}\hbox to\headwidth{\color{cbteal!45}\leaders\hrule height 1.2pt\hfill}}



\fancyhead[L]{\small\color{black!55}Year 5 Maths\sep Unit 5\sep Below Zero in Words}

% A number line from #1 to #2 with every whole number labelled, zero bold.
\newcommand{\numline}[2]{%
  \begin{tikzpicture}[x=0.92cm,y=1cm]
    \draw[line width=1.1pt,cbteal,{Latex[length=3mm]}-{Latex[length=3mm]}] (#1-0.7,0) -- (#2+0.7,0);
    \foreach \n in {#1,...,#2}{
      \draw[line width=0.9pt,cbteal] (\n,-0.18) -- (\n,0.18);
      \node[below=4pt,font=\small] at (\n,-0.18) {$\n$};
    }
    \draw[line width=2pt,cborange] (0,-0.28) -- (0,0.28);
  \end{tikzpicture}}

\begin{document}
\thispagestyle{fancy}

\begin{tcolorbox}[enhanced, colback=cbpurple!7, colframe=cbpurple!7, arc=7pt,
                  borderline west={6pt}{0pt}{cbpurple},
                  left=18pt, right=14pt, top=14pt, bottom=14pt]
  {\color{cbpurple}\bfseries YEAR 5 MATHS\sep UNIT 5\sep WORKSHEET}\\[6pt]
  {\color{cbpurple}\Huge\bfseries Below Zero in Words}\\[10pt]
  {\color{black!70}Adding and subtracting across zero: temperatures, floors and money}\\
  {\color{black!70}Mr. David}
\end{tcolorbox}

\vspace{8pt}
\noindent
\begin{tabularx}{\textwidth}{@{}X X@{}}
  \textbf{Name:} \blank[5.6cm] & \textbf{Class:} \blank[5.6cm] \\
  \textbf{Date given:} \blank[4.6cm] & \textbf{Date due:} \blank[5.0cm] \\
\end{tabularx}

\vspace{6pt}
\begin{activity}[How to do this worksheet]
  \begin{itemize}
    \item Write your answers \textbf{on this sheet}.
    \item In a word problem, \key{write the number sentence first}, then the answer.
    \item Stuck? Draw a number line with \key{0 in the middle}.
    \item A new word? Look for the blue \eng{Word help} box.
  \end{itemize}
\end{activity}

\begin{remember}
  \key{Zero is a bridge.} Jump to 0 first. Then jump the rest.\\[3pt]
  $3-8$: \quad 3 to 0 is 3 jumps, then 5 more jumps. \quad $3-8=-5$
\end{remember}

% ===========================================================================
\hwsection{Section A\sep Number lines}

\qhead{A1}{Count on and count back}
\noindent Write the missing numbers. The steps are always the same size.

\vspace{4pt}
\begin{enumerate}[label=(\alph*), itemsep=8pt]
  \item $9,\ 6,\ 3,\ \blank[1.4cm],\ \blank[1.4cm],\ \blank[1.4cm]$
  \item $-11,\ -7,\ -3,\ \blank[1.4cm],\ \blank[1.4cm],\ \blank[1.4cm]$
  \item $14,\ 9,\ 4,\ \blank[1.4cm],\ \blank[1.4cm],\ \blank[1.4cm]$
  \item $\blank[1.4cm],\ -15,\ -5,\ 5,\ \blank[1.4cm],\ \blank[1.4cm]$
\end{enumerate}

\qhead{A2}{Show the jumps}
\noindent Draw the jumps on the number line. Then write the answer.

\vspace{6pt}
\noindent\textbf{(a)} $3-7=\blank[2cm]$

\vspace{4pt}
\noindent\resizebox{0.92\textwidth}{!}{\numline{-6}{6}}

\vspace{14pt}
\noindent\textbf{(b)} $-5+9=\blank[2cm]$

\vspace{4pt}
\noindent\resizebox{0.92\textwidth}{!}{\numline{-6}{6}}

\qhead{A3}{Work them out}
\noindent The numbers get bigger as you go. Use zero as a bridge.

\vspace{8pt}
\begin{multicols}{3}
\begin{enumerate}[label=(\alph*), itemsep=14pt]
  \item $2-6 = \blank[1.5cm]$
  \item $1-9 = \blank[1.5cm]$
  \item $-4+7 = \blank[1.5cm]$
  \item $-8+3 = \blank[1.5cm]$
  \item $5-12 = \blank[1.5cm]$
  \item $-10+16 = \blank[1.5cm]$
  \item $20-35 = \blank[1.5cm]$
  \item $-25+40 = \blank[1.5cm]$
  \item $7-100 = \blank[1.5cm]$
  \item $-30+120 = \blank[1.5cm]$
  \item $15-204 = \blank[1.5cm]$
  \item $-60+1005 = \blank[1.2cm]$
\end{enumerate}
\end{multicols}

% ===========================================================================
\hwsection{Section B\sep The English of the question}

\begin{wordhelp}[Word help --- these words tell you to add or subtract]
  \begin{tabularx}{\linewidth}{@{}lX@{}}
    \eng{rises}, \eng{goes up}, \eng{gets warmer} & add $(+)$ \\
    \eng{falls}, \eng{drops}, \eng{goes down}, \eng{colder} & subtract $(-)$ \\
    \eng{deeper} & further down: subtract $(-)$ \\
    \eng{ground floor} (G), \eng{B1}, \eng{B2} & floor 0, then $-1$, $-2$ under the ground \\
    \eng{owe} & you must pay money back. Money you owe is \textbf{negative}. \\
    \eng{subtract 5 from 2} & start at 2: $2-5$. The number after \eng{from} comes first. \\
    \eng{the difference between} & the gap between two numbers. It is never negative. \\
  \end{tabularx}
\end{wordhelp}

\qhead{B2}{Number sentence first}
\noindent Write the number sentence. Then write the answer.

\vspace{4pt}
\begin{enumerate}[label=\arabic*., itemsep=9pt]
  \item The temperature is $2\dC$. It falls by $9\dC$. \\[4pt]
        Number sentence: \blank[5.6cm] \quad Answer: \blank[2.6cm]
  \item The lift is on floor 3. It goes down 7 floors. \\[4pt]
        Number sentence: \blank[5.6cm] \quad Floor: \blank[2.6cm]
  \item The lift is on floor B5. It goes up 8 floors. \\[4pt]
        Number sentence: \blank[5.6cm] \quad Floor: \blank[2.6cm]
  \item Mr. David has 15\,000 dong. He buys a book for 40\,000 dong. He borrows the
        rest. How much does he owe? \\[4pt]
        Number sentence: \blank[5.6cm] \quad Answer: \blank[2.6cm]
  \item Subtract 12 from 5. \\[4pt]
        Number sentence: \blank[5.6cm] \quad Answer: \blank[2.6cm]
\end{enumerate}

\qhead{B3}{The difference between}
\noindent Count the gap. A number line helps.

\vspace{4pt}
\begin{enumerate}[label=(\alph*), itemsep=10pt]
  \item The difference between $3$ and $10$. \hfill \blank[2.6cm]
  \item The difference between $-2$ and $6$. \hfill \blank[2.6cm]
  \item The difference between $-9$ and $4$. \hfill \blank[2.6cm]
  \item How much warmer is $8\dC$ than $-5\dC$? \hfill \blank[2.6cm]
  \item How many floors is it from floor B3 up to floor 6? \hfill \blank[2.6cm]
\end{enumerate}

% ===========================================================================
\hwsection{Section C\sep Read the table}

\qhead{C1}{One morning in January}
\noindent The table shows the temperature at 7 a.m. in five places on the same day.
{\small\color{black!60}(The numbers are made up for this worksheet.)}

\vspace{8pt}
\begin{center}
\begin{tabular}{|p{5.2cm}|p{3.4cm}|}
  \hline
  \rule{0pt}{1.5em}\textbf{Place} & \textbf{Temperature} \\ \hline
  \rule{0pt}{1.5em}M\d{\^o}c Ch\^au, Vietnam & $4\dC$ \\ \hline
  \rule{0pt}{1.5em}H\d{a} Long, Vietnam & $13\dC$ \\ \hline
  \rule{0pt}{1.5em}Beijing, China & $-6\dC$ \\ \hline
  \rule{0pt}{1.5em}Ulaanbaatar, Mongolia & $-24\dC$ \\ \hline
  \rule{0pt}{1.5em}Sydney, Australia & $26\dC$ \\ \hline
\end{tabular}
\end{center}

\begin{enumerate}[label=(\alph*), itemsep=11pt]
  \item Which place is the coldest? \hfill \blank[4.2cm]
  \item Beijing gets $9\dC$ warmer. What is the new temperature? \hfill \blank[2.4cm]
  \item In M\d{\^o}c Ch\^au it falls by $7\dC$. What is the new temperature? \hfill \blank[2.4cm]
  \item The difference between Sydney and Beijing? \hfill \blank[2.4cm]
  \item Write the five temperatures in order, \textbf{coldest first}. \\[6pt]
        \blank[2.2cm],\quad \blank[2.2cm],\quad \blank[2.2cm],\quad \blank[2.2cm],\quad \blank[2.2cm]
\end{enumerate}

\qhead{C2}{Estimate first}
\noindent Use easy numbers to estimate. Then work out the exact answer.

\vspace{6pt}
\noindent
\begin{tabularx}{\textwidth}{|p{3.6cm}|X|X|}
  \hline
  \rule{0pt}{1.6em}\textbf{Calculation} & \textbf{Estimate} & \textbf{Exact answer} \\ \hline
  \rule{0pt}{2.2em}$-49+151$ & & \\ \hline
  \rule{0pt}{2.2em}$302-498$ & & \\ \hline
  \rule{0pt}{2.2em}$-99+1001$ & & \\ \hline
\end{tabularx}

% ===========================================================================
\hwsection{Section D\sep Think about it}

\qhead{D1}{Mr. David's homework}
\noindent Mr. David got all three wrong. Write the correct answer. Then say what he did wrong.

\vspace{4pt}
\begin{enumerate}[label=(\alph*), itemsep=6pt]
  \item $6-10=4$ \hfill Correct answer: \blank[2.2cm]
        \wlines{1}
  \item $-3+8=-11$ \hfill Correct answer: \blank[2.2cm]
        \wlines{1}
  \item The difference between $-5$ and $3$ is $2$. \hfill Correct answer: \blank[2.2cm]
        \wlines{1}
\end{enumerate}

\qhead{D2}{Make up the story}
\noindent Write a word problem for each number sentence. Use a temperature, a lift
or money. Use a word from the Word help box. End with a question.

\vspace{6pt}
\noindent\textbf{(a)} \expr{4-9=-5}
\wlines{2}

\vspace{8pt}
\noindent\textbf{(b)} \expr{-3+10=7}
\wlines{2}

% ===========================================================================
% ANSWER KEY (teacher copy only)
% ===========================================================================
\ifkey
\newpage
\hwsection[cbred]{Answer Key --- Teacher Copy}
\linespread{1.0}\selectfont
\setlist[enumerate]{itemsep=1pt, topsep=2pt, leftmargin=*}
\small

\qhead[cbred]{A1}{Count on and count back}
\noindent (a) $0, -3, -6$ \quad (b) $1, 5, 9$ \quad (c) $-1, -6, -11$ \quad
(d) $-25$ \ldots\ $15, 25$ (steps of 10). A student who writes $-1$ after $3$ in (a)
has stopped seeing the step size at zero: count the jumps aloud with them.

\qhead[cbred]{A2}{Show the jumps}
\noindent (a) $-4$: 3 jumps left to 0, then 4 more. \quad (b) $4$: 5 jumps right to 0,
then 4 more. Check the arrow starts at the first number, not at 0.

\qhead[cbred]{A3}{Work them out}
\noindent (a) $-4$ (b) $-8$ (c) $3$ (d) $-5$ (e) $-7$ (f) $6$ (g) $-15$ (h) $15$
(i) $-93$ (j) $90$ (k) $-189$ (l) $945$. \\
The common wrong answers are the gap with the wrong sign ($7-100=93$) and the
``keep the minus'' sum ($-4+7=-11$, the scheme of work's misconception). Both mean
the student is not starting at the first number.

\qhead[cbred]{B2}{Number sentence first}
\noindent 1. $2-9=-7\dC$ \quad 2. $3-7=-4$, floor B4 \quad
3. $-5+8=3$, floor 3 \quad 4. $15\,000-40\,000=-25\,000$: he owes 25\,000 dong \quad
5. $5-12=-7$. \\
In 5, $12-5=7$ is the English mistake (the number after \textit{from} comes first),
not an arithmetic one: ask the student to read the sentence aloud.

\qhead[cbred]{B3}{The difference between}
\noindent (a) 7 \quad (b) 8 \quad (c) 13 \quad (d) $13\dC$ \quad (e) 9 floors (3 up to
G, then 6). Answers of 4, 5 or 3 in (b)--(e) come from subtracting the digits and
ignoring the minus sign. A difference is never negative: $-13$ is wrong.

\qhead[cbred]{C1}{One morning in January}
\noindent (a) Ulaanbaatar \quad (b) $-6+9=3\dC$ \quad (c) $4-7=-3\dC$ \quad
(d) $26+6=32\dC$ \quad (e) $-24, -6, 4, 13, 26$ ($\dC$). In (e), students who put $-6$
before $-24$ are ordering by the digits; ask which is colder.

\qhead[cbred]{C2}{Estimate first}
\noindent $-50+150=100$, exact $102$. \quad $300-500=-200$, exact $-196$. \quad
$-100+1000=900$, exact $902$. Any sensible rounded estimate is fine; the point is that
the exact answer is close to it.

\qhead[cbred]{D1}{Mr. David's homework}
\noindent (a) $-4$: he found the gap between 6 and 10 instead of starting at 6 and going
down 10. (b) $5$: he added 3 and 8 and kept the minus sign. (c) $8$: he subtracted 3
from 5; the gap crosses zero, $5+3$.

\qhead[cbred]{D2}{Make up the story}
\noindent Check in this order: (1) it \textbf{starts} at the first number (4, or $-3$);
(2) the word matches the sign (falls / goes down / spends for (a), rises / goes up /
earns for (b)); (3) full sentences ending in a question; (4) the answer is $-5$ or 7.
Example (a): ``It is $4\dC$. The temperature falls by $9\dC$. What is the temperature
now?'' Example (b): ``The lift is on B3. It goes up 10 floors. Which floor is it on now?''

\fi
\end{document}
"
% or: pwsh .claude/skills/new-homework/scripts/build-hw.ps1 y5m05

\documentclass[12pt,a4paper]{article}

\newif\ifkey
\ifdefined\TEACHER\keytrue\else
  \keyfalse
\fi

% homework/hw-style.tex
% Shared house style for every Year 7 homework packet. \input this from the
% packet file, right after \documentclass and the answer-key switch.
%
% Do not fork this file per packet. If a packet needs a different look, the
% question is whether every packet should have it.
%
% The choices here are deliberate and were arrived at by printing and looking:
%   * No solid dark banners. Every header is a light tint with a coloured spine
%     on the left, so colour reads clearly without flooding a school printer.
%   * No marks and no mark totals. These are practice, not exam papers.
%   * Lato at 12pt with open leading, plus mathastext so inline maths is Lato
%     too. Without mathastext every "-6 + 4" is Computer Modern serif and the
%     sheet looks half-typeset.
%   * Writing lines are spaced for Year 7 handwriting, not adult margin notes.
%   * Pages fill. \needspace is redefined below so a heading no longer drags
%     the page break up to itself, and every heading is glued to the box under
%     it. Packets built before 2026-09-30 were laid out under the old rules and
%     will reflow if rebuilt -- re-read every page if you rebuild one.
%
% Provides: \wlines \blank \tickbox \sep \qmk-free marks (none) \hwsection
% \qhead \expr, the workspace box, and the remember / wordhelp / activity /
% yourturn coloured boxes. Palette matches the lesson decks exactly.

% -- Packages ----------------------------------------------------------------
\usepackage[T1]{fontenc}
\usepackage[utf8]{inputenc}
\usepackage[default]{lato}
\usepackage[a4paper,top=1.7cm,bottom=1.7cm,left=2.0cm,right=2.0cm,headheight=15pt]{geometry}
\usepackage{amsmath,amssymb}
\usepackage{xcolor}
\usepackage[most]{tcolorbox}
\usepackage{enumitem}
\usepackage{tikz}
\usetikzlibrary{arrows.meta}
\usepackage{array}
\usepackage{tabularx}
\usepackage{fancyhdr}
\usepackage{multicol}
\usepackage{needspace}
\usepackage{graphicx}

% Make maths use Lato too. Without this, every "-6 + 4" on the page is set in
% Computer Modern serif and the sheet looks half-typeset. Must come last.
\usepackage[italic]{mathastext}

% -- \needspace, without the hole it used to leave ----------------------------
% needspace.sty writes "\vskip 0pt plus X \penalty-100" in front of every
% heading. A page that ends anywhere else is underfull with no stretch, so TeX
% scores it b=10000; a page that ends AT a heading gets that X of stretch and a
% finite badness, and wins. So any heading in roughly the bottom half of a page
% claimed the break, and the question under it moved over whole, leaving a hole
% the size of the question. \tracingpages showed it exactly (b=4531 at the
% heading, b=10000 at every later break) while building homework/q1rev, and it
% is where the half-empty pages in HW 1-8 came from.
% This keeps only the \penalty9999, which still refuses to put a heading in the
% last X of a page. Keeping a heading WITH the box under it is \nobreak's job
% now -- see \hwsection and \qhead below.
\makeatletter
\renewcommand{\needspace}[1]{\begingroup\setlength{\dimen@}{#1}%
  \vskip\dimen@\penalty9999\vskip-\dimen@\endgroup}
\makeatother

\linespread{1.06}\selectfont

% -- The Learner's Book palette (same hex values as the lesson decks) ---------
\definecolor{cbteal}{HTML}{0087A8}
\definecolor{cbpurple}{HTML}{5C2483}
\definecolor{cborange}{HTML}{C25E12}
\definecolor{cbgreen}{HTML}{4A8B23}
\definecolor{cbred}{HTML}{C8102E}
\definecolor{cbblue}{HTML}{1A5FA8}
\definecolor{rulegray}{HTML}{9AA5AD}

% -- Answer lines and blanks -------------------------------------------------
% \wlines{n} draws n ruled writing lines, spaced wide enough for Year 7
% handwriting rather than for an adult with a fine pen.
\makeatletter
\newcommand{\wlines}[1]{%
  \par\vspace{0.75em}%
  \@tempcnta=0\relax
  \@whilenum\@tempcnta<#1\do{%
    \hbox to \linewidth{\color{rulegray}\leaders\hrule height 0.5pt\hfill}%
    \vspace{1.5em}%
    \advance\@tempcnta by 1\relax}%
  \par\vspace{-0.8em}%
}
\makeatother

% \blank[width] is an inline gap to write a short answer in.
\newcommand{\blank}[1][2.6cm]{\underline{\hspace{#1}}}

% A boxed space for showing working.
\newtcolorbox{workspace}[1][3.4cm]{%
  colback=white, colframe=rulegray, boxrule=0.5pt, arc=5pt,
  height=#1, left=6pt, right=6pt, top=4pt, bottom=4pt,
  before skip=9pt, after skip=11pt}

% An empty tick box.
\newcommand{\tickbox}{\fbox{\rule{0pt}{1.15em}\rule{1.15em}{0pt}}}

% A middot separator.
\newcommand{\sep}{\,\textperiodcentered\,}

% -- Section and question headings -------------------------------------------
% Light tint plus a fat coloured spine on the left. No solid dark banner: the
% colour reads clearly but barely costs the printer anything.
% needspace keeps a heading out of the last few lines of a page. The \nobreak
% after it keeps it with whatever comes next: a remember or word-help box that
% does not fit takes its heading over with it, instead of leaving the heading
% alone at the foot of the page. (The old needspace did that by accident, and
% left a hole every time; see above.)
\newcommand{\hwsection}[2][cbteal]{%
  \needspace{8\baselineskip}%
  \par\vspace{13pt}%
  \noindent
  \begin{tcolorbox}[enhanced, nobeforeafter, width=\textwidth,
      colback=#1!7, colframe=#1!7, arc=5pt,
      borderline west={5pt}{0pt}{#1},
      left=14pt, right=10pt, top=7pt, bottom=7pt]
    {\large\bfseries\color{#1}#2}
  \end{tcolorbox}%
  \par\nobreak\vspace{9pt}}

% \qhead[colour]{A2}{Moving along the line} -- a question heading.
\newcommand{\qhead}[3][cbteal]{%
  \needspace{6\baselineskip}%
  \par\vspace{11pt}%
  \noindent{\large\bfseries\color{#1}#2}\quad{\large\bfseries #3}\par\nobreak\vspace{4pt}}

% -- Coloured boxes ----------------------------------------------------------
% Shared look: rounded, light fill, coloured rule, coloured title text on a
% pale title strip rather than a dark one.
% NOT breakable. A breakable box that cannot fit its first lines forces its own
% page break, and \nobreak cannot stop a forced break -- so the heading above
% it was left alone at the foot of the page (tested: it happens with \nobreak
% in place). Unbreakable, the page builder moves heading and box together.
% These boxes are a few lines each; none needs to split.
\tcbset{hwbox/.style={
  enhanced, arc=6pt, boxrule=1pt,
  left=11pt, right=11pt, top=8pt, bottom=8pt,
  fonttitle=\bfseries, attach boxed title to top left={xshift=12pt, yshift=-8pt},
  boxed title style={arc=4pt, boxrule=1pt, left=7pt, right=7pt, top=2pt, bottom=2pt},
  top=13pt}}

% Orange: things to remember. Matches the deck's "Write This Down".
\newtcolorbox{remember}[1][Remember from the lesson]{hwbox,
  colback=cborange!6, colframe=cborange, coltitle=cborange,
  colbacktitle=white, title={#1}}

% Blue: English help. The barrier in this class is the wording, not the maths.
\newtcolorbox{wordhelp}[1][Word help]{hwbox,
  colback=cbblue!6, colframe=cbblue, coltitle=cbblue,
  colbacktitle=white, title={#1}}

% Purple: an activity or instruction.
\newtcolorbox{activity}[1][How to do this homework]{hwbox,
  colback=cbpurple!6, colframe=cbpurple, coltitle=cbpurple,
  colbacktitle=white, title={#1}}

% Green: your turn / a friendly nudge.
\newtcolorbox{yourturn}[1][Your turn]{hwbox,
  colback=cbgreen!7, colframe=cbgreen, coltitle=cbgreen!75!black,
  colbacktitle=white, title={#1}}

% -- Shorthands --------------------------------------------------------------
\newcommand{\dC}{\,^{\circ}\mathrm{C}}          % use inside maths: $-8\dC$
\newcommand{\key}[1]{{\color{cborange}\textbf{#1}}}
\newcommand{\eng}[1]{{\color{cbblue}\textbf{#1}}}

\setlist[enumerate]{itemsep=8pt, topsep=5pt, leftmargin=*}
\setlist[itemize]{itemsep=4pt, topsep=4pt, leftmargin=*}
\renewcommand{\arraystretch}{1.45}

% -- An expression the student is asked to work from -------------------------
\newcommand{\expr}[1]{%
  \tcbox[on line, colback=cbgreen!10, colframe=cbgreen!55, boxrule=1pt, arc=4pt,
         left=9pt, right=9pt, top=4pt, bottom=4pt]{\large\bfseries $#1$}}

% -- Running head and foot ---------------------------------------------------
% The packet overrides \fancyhead[L] with its own title.
\pagestyle{fancy}
\fancyhf{}
\fancyhead[L]{\small\color{black!55}Year 7\sep Homework}
\fancyhead[R]{\small\color{black!55}Name: \blank[3.4cm]}
\fancyfoot[C]{\small\color{black!55}Page \thepage}
\renewcommand{\headrulewidth}{0pt}
\renewcommand{\headrule}{\vspace{2pt}\hbox to\headwidth{\color{cbteal!45}\leaders\hrule height 1.2pt\hfill}}



\fancyhead[L]{\small\color{black!55}Year 5 Maths\sep Unit 5\sep Below Zero in Words}

% A number line from #1 to #2 with every whole number labelled, zero bold.
\newcommand{\numline}[2]{%
  \begin{tikzpicture}[x=0.92cm,y=1cm]
    \draw[line width=1.1pt,cbteal,{Latex[length=3mm]}-{Latex[length=3mm]}] (#1-0.7,0) -- (#2+0.7,0);
    \foreach \n in {#1,...,#2}{
      \draw[line width=0.9pt,cbteal] (\n,-0.18) -- (\n,0.18);
      \node[below=4pt,font=\small] at (\n,-0.18) {$\n$};
    }
    \draw[line width=2pt,cborange] (0,-0.28) -- (0,0.28);
  \end{tikzpicture}}

\begin{document}
\thispagestyle{fancy}

\begin{tcolorbox}[enhanced, colback=cbpurple!7, colframe=cbpurple!7, arc=7pt,
                  borderline west={6pt}{0pt}{cbpurple},
                  left=18pt, right=14pt, top=14pt, bottom=14pt]
  {\color{cbpurple}\bfseries YEAR 5 MATHS\sep UNIT 5\sep WORKSHEET}\\[6pt]
  {\color{cbpurple}\Huge\bfseries Below Zero in Words}\\[10pt]
  {\color{black!70}Adding and subtracting across zero: temperatures, floors and money}\\
  {\color{black!70}Mr. David}
\end{tcolorbox}

\vspace{8pt}
\noindent
\begin{tabularx}{\textwidth}{@{}X X@{}}
  \textbf{Name:} \blank[5.6cm] & \textbf{Class:} \blank[5.6cm] \\
  \textbf{Date given:} \blank[4.6cm] & \textbf{Date due:} \blank[5.0cm] \\
\end{tabularx}

\vspace{6pt}
\begin{activity}[How to do this worksheet]
  \begin{itemize}
    \item Write your answers \textbf{on this sheet}.
    \item In a word problem, \key{write the number sentence first}, then the answer.
    \item Stuck? Draw a number line with \key{0 in the middle}.
    \item A new word? Look for the blue \eng{Word help} box.
  \end{itemize}
\end{activity}

\begin{remember}
  \key{Zero is a bridge.} Jump to 0 first. Then jump the rest.\\[3pt]
  $3-8$: \quad 3 to 0 is 3 jumps, then 5 more jumps. \quad $3-8=-5$
\end{remember}

% ===========================================================================
\hwsection{Section A\sep Number lines}

\qhead{A1}{Count on and count back}
\noindent Write the missing numbers. The steps are always the same size.

\vspace{4pt}
\begin{enumerate}[label=(\alph*), itemsep=8pt]
  \item $9,\ 6,\ 3,\ \blank[1.4cm],\ \blank[1.4cm],\ \blank[1.4cm]$
  \item $-11,\ -7,\ -3,\ \blank[1.4cm],\ \blank[1.4cm],\ \blank[1.4cm]$
  \item $14,\ 9,\ 4,\ \blank[1.4cm],\ \blank[1.4cm],\ \blank[1.4cm]$
  \item $\blank[1.4cm],\ -15,\ -5,\ 5,\ \blank[1.4cm],\ \blank[1.4cm]$
\end{enumerate}

\qhead{A2}{Show the jumps}
\noindent Draw the jumps on the number line. Then write the answer.

\vspace{6pt}
\noindent\textbf{(a)} $3-7=\blank[2cm]$

\vspace{4pt}
\noindent\resizebox{0.92\textwidth}{!}{\numline{-6}{6}}

\vspace{14pt}
\noindent\textbf{(b)} $-5+9=\blank[2cm]$

\vspace{4pt}
\noindent\resizebox{0.92\textwidth}{!}{\numline{-6}{6}}

\qhead{A3}{Work them out}
\noindent The numbers get bigger as you go. Use zero as a bridge.

\vspace{8pt}
\begin{multicols}{3}
\begin{enumerate}[label=(\alph*), itemsep=14pt]
  \item $2-6 = \blank[1.5cm]$
  \item $1-9 = \blank[1.5cm]$
  \item $-4+7 = \blank[1.5cm]$
  \item $-8+3 = \blank[1.5cm]$
  \item $5-12 = \blank[1.5cm]$
  \item $-10+16 = \blank[1.5cm]$
  \item $20-35 = \blank[1.5cm]$
  \item $-25+40 = \blank[1.5cm]$
  \item $7-100 = \blank[1.5cm]$
  \item $-30+120 = \blank[1.5cm]$
  \item $15-204 = \blank[1.5cm]$
  \item $-60+1005 = \blank[1.2cm]$
\end{enumerate}
\end{multicols}

% ===========================================================================
\hwsection{Section B\sep The English of the question}

\begin{wordhelp}[Word help --- these words tell you to add or subtract]
  \begin{tabularx}{\linewidth}{@{}lX@{}}
    \eng{rises}, \eng{goes up}, \eng{gets warmer} & add $(+)$ \\
    \eng{falls}, \eng{drops}, \eng{goes down}, \eng{colder} & subtract $(-)$ \\
    \eng{deeper} & further down: subtract $(-)$ \\
    \eng{ground floor} (G), \eng{B1}, \eng{B2} & floor 0, then $-1$, $-2$ under the ground \\
    \eng{owe} & you must pay money back. Money you owe is \textbf{negative}. \\
    \eng{subtract 5 from 2} & start at 2: $2-5$. The number after \eng{from} comes first. \\
    \eng{the difference between} & the gap between two numbers. It is never negative. \\
  \end{tabularx}
\end{wordhelp}

\qhead{B2}{Number sentence first}
\noindent Write the number sentence. Then write the answer.

\vspace{4pt}
\begin{enumerate}[label=\arabic*., itemsep=9pt]
  \item The temperature is $2\dC$. It falls by $9\dC$. \\[4pt]
        Number sentence: \blank[5.6cm] \quad Answer: \blank[2.6cm]
  \item The lift is on floor 3. It goes down 7 floors. \\[4pt]
        Number sentence: \blank[5.6cm] \quad Floor: \blank[2.6cm]
  \item The lift is on floor B5. It goes up 8 floors. \\[4pt]
        Number sentence: \blank[5.6cm] \quad Floor: \blank[2.6cm]
  \item Mr. David has 15\,000 dong. He buys a book for 40\,000 dong. He borrows the
        rest. How much does he owe? \\[4pt]
        Number sentence: \blank[5.6cm] \quad Answer: \blank[2.6cm]
  \item Subtract 12 from 5. \\[4pt]
        Number sentence: \blank[5.6cm] \quad Answer: \blank[2.6cm]
\end{enumerate}

\qhead{B3}{The difference between}
\noindent Count the gap. A number line helps.

\vspace{4pt}
\begin{enumerate}[label=(\alph*), itemsep=10pt]
  \item The difference between $3$ and $10$. \hfill \blank[2.6cm]
  \item The difference between $-2$ and $6$. \hfill \blank[2.6cm]
  \item The difference between $-9$ and $4$. \hfill \blank[2.6cm]
  \item How much warmer is $8\dC$ than $-5\dC$? \hfill \blank[2.6cm]
  \item How many floors is it from floor B3 up to floor 6? \hfill \blank[2.6cm]
\end{enumerate}

% ===========================================================================
\hwsection{Section C\sep Read the table}

\qhead{C1}{One morning in January}
\noindent The table shows the temperature at 7 a.m. in five places on the same day.
{\small\color{black!60}(The numbers are made up for this worksheet.)}

\vspace{8pt}
\begin{center}
\begin{tabular}{|p{5.2cm}|p{3.4cm}|}
  \hline
  \rule{0pt}{1.5em}\textbf{Place} & \textbf{Temperature} \\ \hline
  \rule{0pt}{1.5em}M\d{\^o}c Ch\^au, Vietnam & $4\dC$ \\ \hline
  \rule{0pt}{1.5em}H\d{a} Long, Vietnam & $13\dC$ \\ \hline
  \rule{0pt}{1.5em}Beijing, China & $-6\dC$ \\ \hline
  \rule{0pt}{1.5em}Ulaanbaatar, Mongolia & $-24\dC$ \\ \hline
  \rule{0pt}{1.5em}Sydney, Australia & $26\dC$ \\ \hline
\end{tabular}
\end{center}

\begin{enumerate}[label=(\alph*), itemsep=11pt]
  \item Which place is the coldest? \hfill \blank[4.2cm]
  \item Beijing gets $9\dC$ warmer. What is the new temperature? \hfill \blank[2.4cm]
  \item In M\d{\^o}c Ch\^au it falls by $7\dC$. What is the new temperature? \hfill \blank[2.4cm]
  \item The difference between Sydney and Beijing? \hfill \blank[2.4cm]
  \item Write the five temperatures in order, \textbf{coldest first}. \\[6pt]
        \blank[2.2cm],\quad \blank[2.2cm],\quad \blank[2.2cm],\quad \blank[2.2cm],\quad \blank[2.2cm]
\end{enumerate}

\qhead{C2}{Estimate first}
\noindent Use easy numbers to estimate. Then work out the exact answer.

\vspace{6pt}
\noindent
\begin{tabularx}{\textwidth}{|p{3.6cm}|X|X|}
  \hline
  \rule{0pt}{1.6em}\textbf{Calculation} & \textbf{Estimate} & \textbf{Exact answer} \\ \hline
  \rule{0pt}{2.2em}$-49+151$ & & \\ \hline
  \rule{0pt}{2.2em}$302-498$ & & \\ \hline
  \rule{0pt}{2.2em}$-99+1001$ & & \\ \hline
\end{tabularx}

% ===========================================================================
\hwsection{Section D\sep Think about it}

\qhead{D1}{Mr. David's homework}
\noindent Mr. David got all three wrong. Write the correct answer. Then say what he did wrong.

\vspace{4pt}
\begin{enumerate}[label=(\alph*), itemsep=6pt]
  \item $6-10=4$ \hfill Correct answer: \blank[2.2cm]
        \wlines{1}
  \item $-3+8=-11$ \hfill Correct answer: \blank[2.2cm]
        \wlines{1}
  \item The difference between $-5$ and $3$ is $2$. \hfill Correct answer: \blank[2.2cm]
        \wlines{1}
\end{enumerate}

\qhead{D2}{Make up the story}
\noindent Write a word problem for each number sentence. Use a temperature, a lift
or money. Use a word from the Word help box. End with a question.

\vspace{6pt}
\noindent\textbf{(a)} \expr{4-9=-5}
\wlines{2}

\vspace{8pt}
\noindent\textbf{(b)} \expr{-3+10=7}
\wlines{2}

% ===========================================================================
% ANSWER KEY (teacher copy only)
% ===========================================================================
\ifkey
\newpage
\hwsection[cbred]{Answer Key --- Teacher Copy}
\linespread{1.0}\selectfont
\setlist[enumerate]{itemsep=1pt, topsep=2pt, leftmargin=*}
\small

\qhead[cbred]{A1}{Count on and count back}
\noindent (a) $0, -3, -6$ \quad (b) $1, 5, 9$ \quad (c) $-1, -6, -11$ \quad
(d) $-25$ \ldots\ $15, 25$ (steps of 10). A student who writes $-1$ after $3$ in (a)
has stopped seeing the step size at zero: count the jumps aloud with them.

\qhead[cbred]{A2}{Show the jumps}
\noindent (a) $-4$: 3 jumps left to 0, then 4 more. \quad (b) $4$: 5 jumps right to 0,
then 4 more. Check the arrow starts at the first number, not at 0.

\qhead[cbred]{A3}{Work them out}
\noindent (a) $-4$ (b) $-8$ (c) $3$ (d) $-5$ (e) $-7$ (f) $6$ (g) $-15$ (h) $15$
(i) $-93$ (j) $90$ (k) $-189$ (l) $945$. \\
The common wrong answers are the gap with the wrong sign ($7-100=93$) and the
``keep the minus'' sum ($-4+7=-11$, the scheme of work's misconception). Both mean
the student is not starting at the first number.

\qhead[cbred]{B2}{Number sentence first}
\noindent 1. $2-9=-7\dC$ \quad 2. $3-7=-4$, floor B4 \quad
3. $-5+8=3$, floor 3 \quad 4. $15\,000-40\,000=-25\,000$: he owes 25\,000 dong \quad
5. $5-12=-7$. \\
In 5, $12-5=7$ is the English mistake (the number after \textit{from} comes first),
not an arithmetic one: ask the student to read the sentence aloud.

\qhead[cbred]{B3}{The difference between}
\noindent (a) 7 \quad (b) 8 \quad (c) 13 \quad (d) $13\dC$ \quad (e) 9 floors (3 up to
G, then 6). Answers of 4, 5 or 3 in (b)--(e) come from subtracting the digits and
ignoring the minus sign. A difference is never negative: $-13$ is wrong.

\qhead[cbred]{C1}{One morning in January}
\noindent (a) Ulaanbaatar \quad (b) $-6+9=3\dC$ \quad (c) $4-7=-3\dC$ \quad
(d) $26+6=32\dC$ \quad (e) $-24, -6, 4, 13, 26$ ($\dC$). In (e), students who put $-6$
before $-24$ are ordering by the digits; ask which is colder.

\qhead[cbred]{C2}{Estimate first}
\noindent $-50+150=100$, exact $102$. \quad $300-500=-200$, exact $-196$. \quad
$-100+1000=900$, exact $902$. Any sensible rounded estimate is fine; the point is that
the exact answer is close to it.

\qhead[cbred]{D1}{Mr. David's homework}
\noindent (a) $-4$: he found the gap between 6 and 10 instead of starting at 6 and going
down 10. (b) $5$: he added 3 and 8 and kept the minus sign. (c) $8$: he subtracted 3
from 5; the gap crosses zero, $5+3$.

\qhead[cbred]{D2}{Make up the story}
\noindent Check in this order: (1) it \textbf{starts} at the first number (4, or $-3$);
(2) the word matches the sign (falls / goes down / spends for (a), rises / goes up /
earns for (b)); (3) full sentences ending in a question; (4) the answer is $-5$ or 7.
Example (a): ``It is $4\dC$. The temperature falls by $9\dC$. What is the temperature
now?'' Example (b): ``The lift is on B3. It goes up 10 floors. Which floor is it on now?''

\fi
\end{document}
